\chapter{Multiple Integrals}
\section{Double Integrals}
In Chapter 02, we define the integral of a general function \(f: [a,b] \to \R\) by approximating \(f\) by step functions. In this section, we follow the same idea to define the double integral of a function \(f: Q \to \R\) where \(Q\) is a rectangle in \(\R^2\).
\begin{definitionenv}
    Given two closed intervals \([a,b]\) and \([c,d]\) in \(\R\), we call
    \[
        Q := [a,b]\times[c,d]
        = \{(x,y) \in \R^2 : a \leq x \leq b,\ c \leq y \leq d\}
    \]
    a \textbf{rectangle} in \(\R^2\).

    \begin{enumerate}
        \item Suppose \(P_1 = \{x_0, x_1, \dots, x_n\}\) and \(P_2 = \{y_0, y_1, \dots, y_m\}\) are partitions of \([a,b]\) and \([c,d]\), respectively, with
        \[
            a = x_0 < x_1 < \dots < x_n = b,
            \qquad
            c = y_0 < y_1 < \dots < y_m = d.
        \]
        Then we call
        \[
            P := P_1 \times P_2
            =
            \{[x_{i-1},x_i]\times[y_{j-1},y_j] : 1 \leq i \leq n,\ 1 \leq j \leq m\}
        \]
        a \textbf{partition} of \(Q\), which decomposes \(Q\) into \(nm\) \textbf{subrectangles}.

        \item The Cartesian product \((x_{i-1},x_i)\times(y_{j-1},y_j)\), where \(1 \leq i \leq n\) and \(1 \leq j \leq m\), is called an \textbf{open subrectangle} of \(P\), or of \(Q\).

        \item Another partition \(\widetilde P\) is said to be \textbf{finer than} \(P\), or a \textbf{refinement} of \(P\), if \(P \subseteq \widetilde P\), i.e., all the points in \(P\) are also in \(\widetilde P\).

        \item A scalar field \(f\colon Q \to \R\) is called a \textbf{step function} if there exists a partition \(P\) of \(Q\) such that \(f\) is constant on each open subrectangle of \(P\).
    \end{enumerate}
\end{definitionenv}
\begin{remarkenv}
    For a rectangle $Q$ in $\R^2$, the set of all step functions defined on $Q$ is a vector space over $\R$: if $f$ and $g$ are step functions defined on $Q$ with respect to partitions $P_f$ and $P_g$, respectively, then $\alpha f + \beta g$ is a step function defined on $Q$ with respect to the common refinement $P_f \cup P_g$ of $P_f$ and $P_g$ for any $\alpha, \beta \in \R$.
\end{remarkenv}
\begin{definitionenv}
    Suppose $Q=[a,b]\times[c,d]$ with a partition $P=P_1 \times P_2$ where $P_1 = \{x_0, x_1, \dots, x_n\}$ and $P_2 = \{y_0, y_1, \dots, y_m\}$. Let $f$ be a step function defined on $Q$ that takes constant value $c_{ij}$ on the open subinterval $(x_{i-1},x_i)\times(y_{j-1},y_j)$ for $1 \leq i \leq n$ and $1 \leq j \leq m$. Then we define the \textbf{double integral} of $f$ over $Q$ by
    \[
        \iint_Q f(x,y) \dd x \dd y := \sum_{i=1}^n \sum_{j=1}^m c_{ij} \cdot (x_i - x_{i-1})(y_j - y_{j-1}).
    \]
    Note if we replace $P$ with a finer partition $\tilde P$, then the value of the double integral does not change.
\end{definitionenv}
\begin{notationenv}
    We can also write the double integral of $f$ over $Q$ as $\displaystyle \iint_Q f$.
\end{notationenv}
\begin{remarkenv}
    Since $f(x,y) = c_{ij}$ on $(x_{i-1},x_i) \times (y_{j-1},y_j)$, if we let $Q_{ij} := [x_{i-1},x_i]\times[y_{j-1},y_j]$, then we have
    \[
    \begin{aligned}
        \iint_{Q_{ij}} f(x,y) \dd x \dd y &= c_{ij} \cdot (x_i - x_{i-1})(y_j - y_{j-1}) \\ &= f(x,y) \int_{x_{i-1}}^{x_i} \dd x \int_{y_{j-1}}^{y_j} \dd y\\ &= \left( \int_{x_{i-1}}^{x_i} f(x,y) \dd x \right) \left( \int_{y_{j-1}}^{y_j} \dd y  \right)\\ &= \int_{y_{j-1}}^{y_j} \left[ \int_{x_{i-1}}^{x_i} f(x,y) \dd x \right] \dd y.
    \end{aligned}
    \]
    Analogously, we also have
    \[
    \iint_{Q_{ij}} f(x,y) \dd x \dd y = \int_{x_{i-1}}^{x_i} \left[ \int_{y_{j-1}}^{y_j} f(x,y) \dd y \right] \dd x.
    \]
    Summing up all $Q_{ij}$, we get
    \[
    \begin{aligned}
        \iint_Q f(x,y) \dd x \dd y &= \sum_{i=1}^n \sum_{j=1}^m \iint_{Q_{ij}} f(x,y) \dd x \dd y\\ &= \sum_{i=1}^m \int_{y_{j-1}}^{y_j} \left[ \sum_{i=1}^n \int_{x_{i-1}}^{x_i} f(x,y) \dd x \right] \dd y\\ &= \int_c^d \left[ \int_a^b f(x,y) \dd x \right] \dd y\\ &= \int_a^b \left[ \int_c^d f(x,y) \dd y \right] \dd x.
    \end{aligned}
    \]
\end{remarkenv}
\begin{theoremenv}
    Suppose $s$ and $t$ are step functionson rectangle $Q\subseteq \R^2$. Then we have the following properties:
    \begin{enumerate}[label=(\arabic*)]
        \item Linearity: for any $\alpha, \beta \in \R$,
        $$ \iint_Q (\alpha s + \beta t)(x,y) \dd x \dd y = \alpha \iint_Q s(x,y) \dd x \dd y + \beta \iint_Q t(x,y) \dd x \dd y, $$
        \item Additivity: if $Q$ can be subdivided to rectangles $Q_1$ and $Q_2$, then
        $$ \iint_Q s(x,y) \dd x \dd y = \iint_{Q_1} s(x,y) \dd x \dd y + \iint_{Q_2} s(x,y) \dd x \dd y, $$
        \item Comparison: if $s(x,y) \leq t(x,y)$ for all $(x,y) \in Q$, then
        $$ \iint_Q s(x,y) \dd x \dd y \leq \iint_Q t(x,y) \dd x \dd y. $$
    \end{enumerate}
\end{theoremenv}
\begin{definitionenv}
    Let $f$ be a scalar field defined and bounded on a rectangle $Q$. If there exists a unique number $I\in \R$ such that
    $$
    \iint_Q s(x,y) \dd x \dd y \leq I \leq \iint_Q t(x,y) \dd x \dd y
    $$
    for every pair of step functions $s$ and $t$ defined on $Q$ with $s(x,y) \leq f(x,y) \leq t(x,y)$ for all $(x,y) \in Q$, then we call $f$ \textbf{integrable} on $Q$ and call $I$ the \textbf{double integral} of $f$ over $Q$, denoted by
    \[
    I := \iint_Q f(x,y) \dd x \dd y.
    \]
\end{definitionenv}
\begin{remarkenv}
    If $s(x,y) \leq f(x,y)$ for all $(x,y) \in Q$, then we say $s$ is \textbf{below} $f$ and often write $s \leq f$. Similarly, if $f(x,y) \leq t(x,y)$ for all $(x,y) \in Q$, then we say $t$ is \textbf{above} $f$ and often write $f \leq t$.
\end{remarkenv}
\begin{definitionenv}
    Let $f$ be a scalar field defined and bounded on a rectangle $Q\subseteq \R^2$. Let
    $$
    \begin{aligned}
        S&:= \left\{ \iint_Q s(x,y) \dd x \dd y : s \text{ is a step function defined on } Q \text{ with } s \leq f\ \right\},\\ T&:= \left\{ \iint_Q t(x,y) \dd x \dd y : t \text{ is a step function defined on } Q \text{ with } f \leq t\ \right\}.
    \end{aligned}
    $$
    We call $\sup S$ the \textbf{lower integral} of $f$ over $Q$, denoted by $\underline{I}(f)$, and call $\inf T$ the \textbf{upper integral} of $f$ over $Q$, denoted by $\overline{I}(f)$.
\end{definitionenv}
\begin{theoremenv}
    Let $f$ be a scalar field defined and bounded on a rectangle $Q\subseteq \R^2$. Then
    \begin{enumerate}
        \item $\displaystyle \iint_Q s\, \dd x \dd y \leq \underline{I}(f) \leq \overline{I}(f) \leq \iint_Q t\, \dd x \dd y$ for any step functions $s \le f \le t$;
        \item $f$ is integrable on $Q$ if and only if $\underline{I}(f) = \overline{I}(f)$, in which case $\displaystyle \iint_Q f \dd x \dd y = \underline{I}(f) = \overline{I}(f)$.
    \end{enumerate}
\end{theoremenv}
\begin{theoremenv}
    Let $f$ be a scalar field defined, bounded and integrable on a rectangle $Q = [a,b] \times [c,d]$. For every fixed $y \in [c,d]$, suppose $\displaystyle \int_a^b f(x,y) \dd x$ exists and denote it by $A(y)$. If $\displaystyle \int_c^d A(y) \dd y$ exists, then we have
    $$
    \iint_Q f(x,y) \dd x \dd y = \int_c^d A(y) \dd y = \int_c^d \left[ \int_a^b f(x,y) \dd x \right] \dd y.
    $$
\end{theoremenv}
\begin{proofenv}
    For any step functions $s\le f \le t$ on $Q$, we have
    \[
    \int_a^b s(x,y) \dd x \leq \int_a^b f(x,y) \dd x = A(y) \leq \int_a^b t(x,y) \dd x.
    \]
    Then integrating both sides gives
    \[
    \int_c^d \left[ \int_a^b s(x,y) \dd x \right] \dd y \leq \int_c^d A(y) \dd y \leq \int_c^d \left[ \int_a^b t(x,y) \dd x \right] \dd y
    \]
    for any arbitrary step functions $s\le f \le t$ on $Q$. By definition, we have $$\displaystyle \iint_Q f(x,y) \dd x \dd y = \int_c^d A(y) \dd y = \int_c^d \int_a^b f(x,y) \dd x \dd y.$$
\end{proofenv}
\begin{remarkenv}
    Alternatively, we could also evaluate $A(x) \displaystyle = \int_c^d f(x,y) \dd y$  and then calculate $\displaystyle \int_a^b A(x) \dd x$ to get the same result.
\end{remarkenv}
\begin{definitionenv}
    Suppose $f$ is non-negative on a rectangle $Q\subseteq \R^2$, the set
    $$
    \mathcal{S} := \{(x,y,z) \in \R^3 : (x,y) \in Q, 0 \leq z \leq f(x,y)\}
    $$
    is called the \textbf{ordinate set} of $f$ over $Q$.
\end{definitionenv}
\begin{remarkenv}
    If $f$ is non-negative and integrable on $Q$, then $\displaystyle \iint_Q f(x,y) \dd x \dd y = \text{vol}(\mathcal{S})$, where $\text{vol}(\mathcal{S})$ is the volume of $\mathcal{S}$. Particularly, for any $y_0 \in [c,d]$, $\displaystyle A(y_0) = \int_a^b f(x,y_0) \dd x$ describes the area of the cross section of $\mathcal{S}$ at $y=y_0$. Moreover, $\displaystyle \iint_Q f(x,y) \dd x \dd y = \int_c^d A(y) \dd y$ can be interpreted that the volume of $\mathcal{S}_f$ is the integral of the cross-sectional area function $A(y)$ from $y=c$ to $y=d$. 
\end{remarkenv}
\begin{theoremenv}[Fubini's Theorem]
    If a scalar field $f$ is continuous on a rectangle $Q = [a,b] \times [c,d]$, then $f$ is integrable on $Q$. Moreover, we have
    $$ \iint_Q f(x,y) \dd x \dd y = \int_c^d \left[ \int_a^b f(x,y) \dd x \right] \dd y = \int_a^b \left[ \int_c^d f(x,y) \dd y \right] \dd x. $$
\end{theoremenv}
\begin{proofenv}[Proof of Integrability]
    Since $f$ is continuous on $Q$, $f$ is also bounded on $Q$ and by Theorem~\ref{thm:small-span}, we know given $\varepsilon > 0$, there is a partition $P$ of $Q$ such that the span of $f$ on each of the closed subrectangles in $P$ is less than $\varepsilon$. For convenience, suppose $P$ decomposes $Q$ into $Q_1,\ldots, Q_N$ subrectangles and the span of $f$ on $Q_k$ is denoted by $M_k(f) - m_k(f)$ and satisfies $M_k(f) - m_k(f) < \varepsilon$ for $k = 1, \ldots, N$. Then we can define a pair of step functions $s$ and $t$ on $Q$ as:
    \begin{itemize}
        \item for $x\in \operatorname{int} Q_k$, let $s(x) := m_k(f)$ and $t(x) := M_k(f)$;
        \item for $x\in \partial Q_k$, let $s(x) := m$ and $t(x) := M$, where $m$ and $M$ are the absolute minimum and maximum values of $f$ on $Q$, respectively.
    \end{itemize}
    Then we have $s \leq f \leq t$ on $Q$ and
    \[
    \iint_Q s = \sum_{k=1}^N m_k(f) \cdot A(Q_k), \qquad \iint_Q t = \sum_{k=1}^N M_k(f) \cdot A(Q_k),
    \]
    where $A(\cdot)$ denotes the area function. Thus, we have
    \[
    0 \le \iint_Q t - \iint_Q s = \sum_{k=1}^N (M_k(f) - m_k(f)) \cdot A(Q_k) < \varepsilon \sum_{k=1}^N A(Q_k) = \varepsilon A(Q).
    \]
    Note
    $$
    0 \le \overline{I}(f) - \underline{I}(f) \leq \iint_Q t - \iint_Q s < \varepsilon A(Q),
    $$
    which implies $\overline{I}(f) = \underline{I}(f)$ since $\varepsilon$ can be arbitrarily small. Therefore, $f$ is integrable on $Q$.
\end{proofenv}
\begin{proofenv}
    For any fixed $y\in [c,d]$, $f(x,y)$ is a continuous function of $x$ on $[a,b]$, so $\displaystyle A(y) := \int_a^b f(x,y) \dd x$ is well-defined. Next, we need to show $A(y)$ is integrable on $[c,d]$. We will do this by showing that $A(y)$ is continuous on $[c,d]$. For any $y, y_0 \in [c,d]$, we have
    \[
    |A(y) - A(y_0)| = \left| \int_a^b f(x,y) \dd x - \int_a^b f(x,y_0) \dd x \right| \le \int_a^b |f(x,y) - f(x,y_0)| \dd x \le (b-a) \cdot \left|f(x_0, y) - f(x_0, y_0) \right|
    \]
    for some $x_0 \in [a,b]$. Thus
    \[
    |A(y) - A(y_0)| \to 0 \quad \text{as } y \to y_0,
    \]
    which implies $A(y)$ is continuous at every $y_0 \in [c,d]$. Then
    $$ \iint_Q f(x,y) \dd x \dd y = \int_c^d \left[ \int_a^b f(x,y) \dd x \right] \dd y. $$
\end{proofenv}
\begin{remarkenv}
    There is some functions $f:\R^2 \to \R$ not continuous on $Q$ but the set of discontinuities of $f$ is ``not too large'', then $f$ still remains integrable on $Q$. 
\end{remarkenv}
\begin{definitionenv}
    Let $B$ be a bounded set in $\R^2$. $B$ is said to \textbf{have content zero} if, for every $\varepsilon >0$, there is a finite collection of rectangles whose union contains $B$ and the sum of whose areas does not exceed $\varepsilon$. Formally, there exists rectangles $R_1, \ldots, R_k$ such that
    $$
    B\subseteq \bigcup_{i=1}^k R_i, \qquad \sum_{i=1}^k A(R_i) < \varepsilon,
    $$
    where $A(\cdot)$ denotes the area function.
\end{definitionenv}
\begin{exampleenv}
    \begin{enumerate}
        \item A set containing finitely many points in $\R^2$ has content zero.
        \item A line segment in $\R^2$ has content zero.
        \item If $A_1, A_2, \ldots, A_n$ all have content zero, then $\displaystyle \bigcup_{i=1}^n A_i$ has content zero.
        \item If $A$ has content zero and $B \subseteq A$, then $B$ has content zero.
    \end{enumerate}
\end{exampleenv}
\begin{theoremenv}
    If a scalar field $f$ is defined and bounded on a rectangle $Q$ and the set of discontinuities of $f$ has content zero, then $f$ is integrable on $Q$.
\end{theoremenv}
\begin{proofenv}
    Let $M>0$ be such that $|f(x,y)| \le M$ for every $(x,y) \in Q$. Let
    $$ D = \{(x,y) \in Q : f \text{ is not continuous at } (x,y)\} $$
    be the set of discontinuities of $f$. Since $D$ has content zero, for every $\varepsilon > 0$, we can pick a partition $P$ of $Q$ such that:
    \begin{itemize}
        \item suppose $\{ Q_1, \ldots, Q_d \}$ denotes the closed subrectangles of $P$ such that
        $$ Q_k \cap D \neq \varnothing,\qquad D \subseteq \bigcup_{k=1}^d Q_k\qquad \text{with } \sum_{k=1}^d A(Q_k) < \varepsilon.$$
        \item suppose $\{Q_{d+1}, \ldots, Q_N\}$ denotes the rest of the closed subrectangles in $P$. Then $f$ is continuous on each $Q_k$ for $k = d+1, \ldots, N$. Furthermore, by Theorem~\ref{thm:small-span}, the span of $f$ on each $Q_k$, denoted by $M_k(f) - m_k(f)$, is less than $\varepsilon$ for $k = d+1, \ldots, N$.
    \end{itemize}
    Then we can define a pair of step functions $s$ and $t$ on $Q$ as
    $$
    s(x) := \begin{cases}
        -M & \text{if } x \in \operatorname{int} Q_k \text{ for } 1 \le k \le d,\\
        m_k(f) & \text{if } x \in \operatorname{int} Q_k \text{ for } d+1 \le k \le N,\\
        -M & \text{if } x \in \partial Q_k \text{ for } 1 \le k \le N,
    \end{cases}\quad
    t(x) := \begin{cases}
        M & \text{if } x \in \operatorname{int} Q_k \text{ for } 1 \le k \le d,\\
        M_k(f) & \text{if } x \in \operatorname{int} Q_k \text{ for } d+1 \le k \le N,\\
        M & \text{if } x \in \partial Q_k \text{ for } 1 \le k \le N.
    \end{cases}
    $$
    Then $s\le f \le t$ on $Q$ and
    $$
    0\le \overline{I}(f) - \underline{I}(f) \le \iint_Q t - \iint_Q s = \sum_{k=1}^d 2M\cdot A(Q_k) + \sum_{k=d+1}^N (M_k(f) - m_k(f)) \cdot A(Q_k) < 2M\varepsilon + \varepsilon A(Q)
    $$
    for any $\varepsilon > 0$. Thus, $\overline{I}(f) = \underline{I}(f)$, which implies $f$ is integrable on $Q$.
\end{proofenv}

\section{Double Integrals over more General Regions}
\begin{definitionenv}
    Suppose $S\subseteq \R^2$ is a bounded region and there is a rectangle $Q=[a,b] \times [c,d]$ such that $S \subseteq Q$. Let $f$ be a scalar field defined and bounded on $S$. We define the \textbf{etended scalar field} $\tilde f$ of $f$ on $Q$ by
    $$
    \tilde f(x,y) := \begin{cases}
    f(x,y) & \text{if } (x,y) \in S,\\
    0 & \text{if } (x,y) \in Q\setminus S.
    \end{cases}
    $$
    If $\tilde f$ is integrable on $Q$, then we say $f$ is \textbf{integrable} on $S$ and define the \textbf{double integral} of $f$ over $S$ by
    $$\iint_S f(x,y) \dd x \dd y := \iint_Q \tilde f(x,y) \dd x \dd y.$$
\end{definitionenv}
\begin{remarkenv}
    We need to show that the definition does not depend on the choice of the rectangle $Q$. Otherwise, the definition is not well-defined. Suppose $Q_1=[a_1,b_1] \times [c_1,d_1]$ and $Q_2=[a_2,b_2] \times [c_2,d_2]$ are two arbitrary rectangles containing $S$. Now define $\tilde f_1$ and $\tilde f_2$ to be the extended scalar fields of $f$ on $Q_1$ and $Q_2$, respectively. Note $S\subseteq Q_1 \cap Q_2$. Thus if $\tilde f_1$ is integrable on $Q_1$, then $\tilde f_1$ is also integrable on $Q_1 \cap Q_2$ and
    $$\iint_{Q_1} \tilde f_1 = \iint_{Q_1 \cap Q_2} \tilde f_1 = \iint_{Q_1 \cap Q_2} \tilde f_2 = \iint_{Q_2} \tilde f_2.$$
    Therefore, the definition of the double integral of $f$ over $S$ is invariant under the choice of the rectangle $Q$.
\end{remarkenv}
\begin{definitionenv}
    We define two types of regions in $\R^2$ as follows:
    \begin{enumerate}[label=(\arabic*)]
        \item a subset $S_I$ of $\R^2$ is called a \textbf{type I region} if there are two continuous functions $\varphi_1, \varphi_2: [a,b] \to \R$ satisfying $\varphi_1(x) \leq \varphi_2(x)$ for all $x \in [a,b]$ such that
        $$ S_I = \{(x,y) \in \R^2 : a \leq x \leq b, \varphi_1(x) \leq y \leq \varphi_2(x)\}; $$
        \item a subset $S_{II}$ of $\R^2$ is called a \textbf{type II region} if there are two continuous functions $\psi_1, \psi_2: [c,d] \to \R$ satisfying $\psi_1(y) \leq \psi_2(y)$ for all $y \in [c,d]$ such that
        $$ S_{II} = \{(x,y) \in \R^2 : c \leq y \leq d, \psi_1(y) \leq x \leq \psi_2(y)\}. $$
    \end{enumerate}
\end{definitionenv}
\begin{remarkenv}
    We will consider integrals $\displaystyle \iint_S f$ for regions $S\subseteq \R^2$ that are of type I, type II, or that can be partitioned into finitely many subregions of type I or type II.
\end{remarkenv}
\begin{theoremenv}
    If $\varphi: [a,b] \to \R$ is a continuous function, then the graph of $\varphi$, i.e., the set $\{(x,y) \in \R^2 : y = \varphi(x), a \leq x \leq b\}$, has content zero in $\R^2$.
\end{theoremenv}
\begin{theoremenv}
    Let $S\subseteq \R^2$ be a region of type I between the graphs of two continuous functions $\varphi_1, \varphi_2: [a,b] \to \R$. Suppose $f:S\to \R$ is a scalar field bounded on $S$ and continuous on $\operatorname{int} S$. Then $f$ is integrable on $S$ and its value is given by
    $$ \iint_S f(x,y) \dd x \dd y = \int_a^b \left[ \int_{\varphi_1(x)}^{\varphi_2(x)} f(x,y) \dd y \right] \dd x. $$
\end{theoremenv}
\begin{corollaryenv}
    Let $S\subseteq \R^2$ be a region of type II between the graphs of two continuous functions $\psi_1, \psi_2: [c,d] \to \R$. Suppose $f:S\to \R$ is a scalar field bounded on $S$ and continuous on $\operatorname{int} S$. Then $f$ is integrable on $S$ and its value is given by
    $$ \iint_S f(x,y) \dd x \dd y = \int_c^d \left[ \int_{\psi_1(y)}^{\psi_2(y)} f(x,y) \dd x \right] \dd y. $$
\end{corollaryenv}
\begin{exampleenv}
    The solid enclosed by the ellipsoid $\displaystyle \frac{x^2}{a^2} + \frac{y^2}{b^2} + \frac{z^2}{c^2} = 1$ is given by the set
    $$
    E_{a,b,c} := \left\{(x,y,z) \in \R^3 : \frac{x^2}{a^2} + \frac{y^2}{b^2} + \frac{z^2}{c^2} \leq 1\right\}.
    $$
    Let
    $$S_{a,b} = \left\{(x,y) \in \R^2 : \frac{x^2}{a^2} + \frac{y^2}{b^2} \leq 1\right\}$$
    and define $f_{a,b,c}: S_{a,b} \to \R$ by
    $$
    f(x,y) = c \sqrt{1 - \frac{x^2}{a^2} - \frac{y^2}{b^2}}.
    $$
    Then $f_{a,b,c}$ is continuous on $S_{a,b}$ and the ordinate set of $f$ over $S_{a,b}$ is exactly $E_{a,b,c}$. Thus, we have
    \[
    \begin{aligned}
        \text{vol}(E_{a,b,c}) &= \iint_{S_{a,b}} \left[ f(x,y) - (-f(x,y)) \right] \dd x \dd y\\
        &= 2 \int_{-a}^a \left[ \int_{-b \sqrt{1-(x/a)^2}}^{b \sqrt{1-(x/a)^2}} c \sqrt{1 - \frac{x^2}{a^2} - \frac{y^2}{b^2}} \dd y \right] \dd x\\
        &= 2 \int_{-a}^a \left[ \int_{-bA}^{bA} c \sqrt{A^2 - \frac{y^2}{b^2}} \dd y \right] \dd x \quad \text{where } A = \sqrt{1 - \frac{x^2}{a^2}}\\
        &= 2 \int_{-a}^a A^2bc \left[ \int_{-\pi/2}^{\pi/2} \cos^2 \theta \dd \theta \right] \dd x \quad \text{where } y = bA \sin \theta,\theta \in [-\pi/2, \pi/2]\\
        &= 2 \int_{-a}^a A^2bc \frac{\pi}{2} \dd x \\
        &= \pi abc \int_{-a}^a \left(1 - \frac{x^2}{a^2}\right) \dd x\\
        &= \frac{4}{3} \pi abc.
    \end{aligned}
    \]
\end{exampleenv}
\section{Green's Theorem}
\begin{definitionenv}
    Suppose $C$ is a curve described by a continuous vector-valued function $\bm \alpha: [a,b] \to \R^n$. Then
    \begin{enumerate}[label=(\arabic*)]
        \item we say $C$ is a \textbf{closed curve} if $\bm \alpha(a) = \bm \alpha(b)$;
        \item we say $C$ is a \textbf{simple closed curve} if $C$ is a closed curve and $\bm \alpha(t_1) \neq \bm \alpha(t_2)$ for any $t_1, t_2 \in [a,b)$ with $t_1 \neq t_2$.
        \item we say $C$ is a \textbf{Jordan curve} if it is a simple closed curve in a plane, i.e., $\alpha : [a,b] \to \R^2$.
        \item The Jordan Curve Theorem states that if $C$ is a Jordan curve, then $\R^2 \setminus C$ is the disjoint union of two connected open subsets. One of these components is bounded and is called the \textbf{inner region}. The other component is unbounded and is called the \textbf{outer region}. Moreover, $C$ is the common boundary of the inner and outer regions.
    \end{enumerate}
\end{definitionenv}
\begin{theoremenv}[Green's Theorem in the Plane]
    Let $P$ and $Q$ be two scalar fields defined and continuously differentiable on an open open set $S\subseteq \R^2$. Let $C$ be a piecewise smooth Jordan curve and let $R$ be the region enclosed by $C$ (containing $C$). If $R \subseteq S$, then we have
    $$ \oint_C P \dd x + Q \dd y = \iint_R \left( \pdv{Q}{x} - \pdv{P}{y} \right) \dd x \dd y. $$
\end{theoremenv}
\begin{exampleenv}
    Let $C$ be a piecewise smooth Jordan curve and let $R$ be the region enclosed by $C$. Then we have
    \[
    A(R) = \iint_R 1 \dd x \dd y = \iint_R \left( \pdv{Q}{x} - \pdv{P}{y} \right) \dd x \dd y = \oint_C P \dd x + Q \dd y = \frac12 \oint_C x \dd y - y \dd x,
    \]
    where we take $Q(x,y) =  \dfrac{x}{2}$ and $P(x,y) = -\dfrac{y}{2}$ to apply Green's Theorem. In particular, if $C$ is described by the parameterization $x = X(t)$ and $y = Y(t)$ for $t \in [a,b]$, then
    \[
    A(R) = \frac12 \int_a^b X(t) Y'(t) - Y(t) X'(t) \dd t = \frac12 \int_a^b \det \begin{pmatrix} X(t) & Y(t)\\ X'(t) & Y'(t) \end{pmatrix} \dd t.
    \]
\end{exampleenv}